\section{Conclusion}
\label{sec:conclusion}

We presented a systematic cross-subject evaluation of three EEG foundation
models for resting-state stress classification, validated by two clinical
reference datasets. Three claims are robust to method choice:

\begin{enumerate}[topsep=0.25em,itemsep=0.1em,leftmargin=1.5em]
  \item Trial-level cross-validation inflates reported balanced accuracy
    by $\sim$21~points relative to subject-level cross-validation, across
    three architecturally distinct foundation models.
  \item Frozen FM features encode subject identity 11--13$\times$ more
    than diagnostic labels, on both stress and Alzheimer's data.
  \item Fine-tuning amplifies label signal only when the underlying
    biological signal is strong enough; weak signals (chronic stress)
    saturate at the same $\sim$0.66 \BA{} achievable with classical
    band-power features and a random forest, and standard mitigations
    (subject-adversarial, LEAD-style auxiliary losses) degrade rather
    than improve this.
\end{enumerate}

We do not interpret this as failure of resting-state EEG for stress
detection. The same data shows that a real but modest stress signature
exists --- and the same diagnosis tells us which research direction will
exceed the ceiling: subject identity should be treated as a reference
frame, not as noise to be removed. Within-subject longitudinal modelling
is, in our reading, the productive next experiment.
